% ============================================================
\subsection{Resolution-Integrated Balance Closure and Hydrodynamic Readout}
\label{sec:flagship-hydrodynamics}
% ============================================================

The hydrodynamic lane enters only after the common selected-cut law has
already been fixed. One begins with exact kinetic balance on observables
preserved by unresolved exchange and then asks what remains visible after the
higher-moment tail has been compressed to the selected endpoint. In that
sense, the primitive object of the section is not the Navier--Stokes system
itself, but the collision-invariant moment hierarchy from which its
recognizable continuum form is read. This follows the standard kinetic-to-
hydrodynamic direction of classical moment theory and projection-based
coarse graining
\cite{ChapmanCowling1970,Cercignani1988,Zwanzig1961,ChorinHaldKupferman2000},
but here the closure coefficients are read from the selected-cut endpoint map
forced by the earlier carrier calculus.

The hydrodynamic lane is therefore not an additional intrinsic persistence
mode. It is a \emph{resolution-integrated effective closure}: only low-order
balances remain visible once unresolved exchange has been integrated over a
resolution family through the tower calculus and compressed into a coarse law.
The exact visible content of this lane begins with the moment hierarchy on the
collision-invariant carrier. Exact moment balance, characteristic endpoint
reduction, and isotropic channel scalarization then determine the endpoint law
class. What is intrinsic here is the underlying exact balance hierarchy
together with its exchange bookkeeping; the recognizable Navier--Stokes form
is the effective constitutive closure read after unresolved kinetic-sector
return has been compressed into endpoint response.

On this lane,
Theorem~\ref{thm:temporal-realization-forbids-primitive-internal-creation}
interprets production terms as balanced exchange with unresolved kinetic
sectors. Collision invariants are precisely the visible observables for which
that exchange source vanishes, so the hydrodynamic law is not presented on the
same footing as the intrinsic reversible phase, wave, gauge, and geometric
sectors.

The relevant exact object on this lane is the balance hierarchy: the selected
cut is the smallest one on which the low-order moment balances are faithfully
visible, and the recognizable hydrodynamic law appears only after closing the
unresolved higher-moment return induced by that hierarchy. The time variable
appearing below is therefore the coordinate expression of the realized
temporal order on the kinetic lane.

We therefore proceed in the same order as the underlying mathematics: first
the exact kinetic hierarchy, then the least carrier on which its visible
balances close, and only then the induced continuum endpoint law.

\subsubsection{Kinetic Realization and Moment Hierarchy}

\begin{definition}[Characteristic kinetic realization]
\label{def:kinetic-transport-branch}\lean{HydrodynamicLane.KineticRealization}
Let \(X\subseteq \mathbb R^d\) be a spatial domain, let \(\Xi\subseteq \mathbb
R^d\) be a velocity space with measure \(d\mu(\xi)\), and let
\[
f : X\times \Xi\times \mathbb R \to \mathbb R
\]
be a resolved state. A \emph{characteristic kinetic realization} consists of a streaming
field
\[
v : X\times \Xi\times \mathbb R \to \mathbb R^d,
\]
a gain kernel \(K(x,t;\xi,\eta)\), and a loss rate \(\sigma(x,t;\xi)\) such
that \(f\) satisfies
\[
\partial_t f(x,\xi,t)
+
\nabla\!\cdot\bigl(v(x,\xi,t)f(x,\xi,t)\bigr)
=
\int_\Xi K(x,t;\xi,\eta)f(x,\eta,t)\,d\mu(\eta)
-\sigma(x,t;\xi)f(x,\xi,t).
\]
\end{definition}

\begin{definition}[Kinetic moment observable]
\label{def:kinetic-moment-observable}\lean{HydrodynamicLane.MomentObservable}
Assume the hypotheses of Definition~\ref{def:kinetic-transport-branch}. Let
\[
\varphi:\Xi\to\mathbb R
\]
be a measurable observable. Its associated moment density, moment flux, and
collision production (equivalently, internal exchange production) are
\[
M_\varphi(x,t)
:=
\int_\Xi \varphi(\xi)f(x,\xi,t)\,d\mu(\xi),
\]
\[
F_\varphi(x,t)
:=
\int_\Xi \varphi(\xi)v(x,\xi,t)f(x,\xi,t)\,d\mu(\xi),
\]
\[
Q_\varphi(x,t)
:=
\int_\Xi\int_\Xi
\bigl(\varphi(\xi)-\varphi(\eta)\bigr)
K(x,t;\xi,\eta)f(x,\eta,t)\,d\mu(\eta)\,d\mu(\xi).
\]
\end{definition}

\begin{theorem}[Kinetic moment hierarchy]
\label{thm:kinetic-moment-hierarchy}\lean{HydrodynamicLane.kineticMomentHierarchy}
Assume the hypotheses of Definition~\ref{def:kinetic-moment-observable}. Then
\[
\partial_t M_\varphi(x,t)
+
\nabla\!\cdot F_\varphi(x,t)
=
Q_\varphi(x,t).
\]
In particular, if \(\varphi\) is a collision invariant in the sense that
\[
\int_\Xi \varphi(\xi)K(x,t;\xi,\eta)\,d\mu(\xi)
=
\varphi(\eta)\sigma(x,t;\eta)
\]
for a.e. \((x,t,\eta)\), then
\[
Q_\varphi(x,t)=0
\]
and hence
\[
\partial_t M_\varphi(x,t)
+
\nabla\!\cdot F_\varphi(x,t)=0.
\]
\end{theorem}

(Proof in Appendix~\ref{sec:appendix-dynamics}.)

On the selected cut,
Theorem~\ref{thm:temporal-realization-forbids-primitive-internal-creation}
interprets \(Q_\varphi\) as the visible exchange source generated by
unresolved kinetic transfer. Collision invariants are
therefore exactly those observables for which the selected visible exchange
source vanishes identically.

\begin{definition}[Hydrodynamic carrier and velocity moments]
\label{def:velocity-channel-kinetic-branch}\lean{HydrodynamicLane.CarrierAndVelocityMoments}
Assume the hypotheses of Definition~\ref{def:kinetic-transport-branch}.
Assume moreover that \(\Xi\subseteq \mathbb R^d\) and that the streaming field
is the native velocity:
\[
v(x,\xi,t)=\xi.
\]
Define the first three observed moments by
\[
U_0(x,t):=\int_\Xi f(x,\xi,t)\,d\mu(\xi),
\]
\[
U_1(x,t):=\int_\Xi \xi\,f(x,\xi,t)\,d\mu(\xi),
\]
\[
U_2(x,t):=\int_\Xi \xi\otimes \xi\,f(x,\xi,t)\,d\mu(\xi).
\]
\[
\mathcal X_{\mathrm{hyd}}
:=
\operatorname{span}\{1,\xi_1,\dots,\xi_d\}.
\]
\end{definition}

\subsubsection{Least-Motion Hydrodynamic Carrier}

\begin{theorem}[Least-motion hydrodynamic observer]
\label{thm:least-motion-hydrodynamic-observer}\lean{HydrodynamicLane.leastMotionObserver}
Assume the hypotheses of Definition~\ref{def:velocity-channel-kinetic-branch}.
Suppose:
\begin{enumerate}[label=(\roman*),itemsep=0.2em]
\item every observable in \(\mathcal X_{\mathrm{hyd}}\) is a collision
      invariant of the realized kinetic law;
\item the carrier \(\mathcal X_{\mathrm{hyd}}\) is faithful,
      closure-admissible, and stable under admissible promotion for the
      realized kinetic family;
\item no proper subcarrier of \(\mathcal X_{\mathrm{hyd}}\) still gives exact
      visible mass-momentum balance together with faithful reconstruction and
      stable promotion at the observed scale;
\item any characteristic observer datum satisfying \emph{(i)}--\emph{(iii)} is
      horizon-preserving equivalent to the one generated by
      \(\mathcal X_{\mathrm{hyd}}\).
\end{enumerate}
Then the carrier generated by \(1,\xi_1,\dots,\xi_d\) is the unique
least-motion characteristic observer for the realized kinetic family.
\end{theorem}

(Proof in Appendix~\ref{sec:appendix-dynamics}.)

\begin{corollary}[Hydrodynamic moment balance]
\label{cor:hydrodynamic-moment-balance}\lean{HydrodynamicLane.momentBalance}
Assume the hypotheses of
Theorem~\ref{thm:least-motion-hydrodynamic-observer}. Then
\[
\partial_t U_0 + \nabla\!\cdot U_1 = 0
\]
and
\[
\partial_t U_1 + \nabla\!\cdot U_2 = 0,
\]
where the divergence of \(U_2\) is taken row by row.
\end{corollary}

(Proof in Appendix~\ref{sec:appendix-dynamics}.)

\subsubsection{Induced Hydrodynamic Law}

At this stage the effective character of the hydrodynamic lane first becomes
explicit. The low-order mass-momentum balances are exact on the least-motion
carrier, but the stress channel still records unresolved higher-moment return.
The next theorem identifies the first selected-cut endpoint law obtained by
reducing that return to isotropic boundary data.

\begin{theorem}[Induced hydrodynamic law on the selected cut]
\label{cor:hydrodynamic-closure-law}\lean{HydrodynamicLane.closureLaw}
Assume the hypotheses of Corollary~\ref{cor:hydrodynamic-moment-balance}.
Assume that \(U_0>0\) and define
\[
w:=U_0^{-1}U_1.
\]
Define the scalar and traceless channels by
\[
s:=\frac{1}{d}\operatorname{tr}\bigl(U_2-U_0\,w\otimes w\bigr),
\qquad
\Theta:=U_2-U_0\,w\otimes w-sI.
\]
Assume moreover that the selected cut satisfies the following characteristic
return properties:
\begin{enumerate}[label=(\roman*),itemsep=0.2em]
\item for each \((x,t)\), with \(\gamma_{x,t}\) the realized characteristic
      through \((x,t)\), there is a rapidly decaying memory family
      \(\mathcal K_{x,t}(s)\) on \(s\le 0\) such that
      \[
      \Theta(x,t)
      =
      \int_{-\infty}^0
      \mathcal K_{x,t}(s)\bigl[\nabla w(\gamma_{x,t}(s),t+s)\bigr]\,ds;
      \]
\item the bulk contribution of this characteristic integral reduces to endpoint
      data by Corollary~\ref{cor:characteristic-boundary-reduction};
\item on the decay scale of \(\mathcal K_{x,t}\), the field \(w\) varies
      slowly enough that the characteristic integral admits a first-order
      endpoint expansion;
\item the induced endpoint map is objective and has no preferred direction.
\end{enumerate}
Then there are unique scalar transport coefficients
\(\mu_{\mathrm{ind}},\zeta_{\mathrm{ind}}\), induced by the selected cut, and a
higher-order remainder \(\mathsf R_{\mathrm{hyd}}\) such that
\[
\Theta
=
-2\mu_{\mathrm{ind}}\,\mathsf D_0(w)
-\zeta_{\mathrm{ind}}(\nabla\!\cdot w)I
+\mathsf R_{\mathrm{hyd}},
\]
where
\[
\mathsf D_0(w)
:=
\frac{1}{2}\bigl(\nabla w+\nabla w^\top\bigr)
-\frac{1}{d}(\nabla\!\cdot w)I.
\]
Consequently
\[
\partial_t U_0 + \nabla\!\cdot(U_0 w)=0
\]
and
\[
\partial_t(U_0 w)
+
\nabla\!\cdot(U_0 w\otimes w)
+
\nabla s
=
\nabla\!\cdot\Bigl(
2\mu_{\mathrm{ind}}\,\mathsf D_0(w)
+
\zeta_{\mathrm{ind}}(\nabla\!\cdot w)I
\Bigr)
-\nabla\!\cdot \mathsf R_{\mathrm{hyd}}.
\]
\end{theorem}

(Proof in Appendix~\ref{sec:appendix-dynamics}.)

\subsubsection{Effective Closure Identification}

The classification point can now be stated cleanly. The exact law on this lane
is the moment hierarchy together with its selected endpoint reduction; the
textbook continuum equation is the effective constitutive readout of that
exact hierarchy. In the standard kinetic literature this endpoint is the
compressible Navier--Stokes law
\cite{ChapmanCowling1970,Cercignani1988,Navier1823,Stokes1845}, and in the
present framework it appears with the same status: as the recognizable
low-order closure induced by unresolved kinetic return, not as an additional
primitive sector law.

\begin{corollary}[Recognizable Navier--Stokes identity]
\label{cor:recognizable-navier-stokes-form}\lean{HydrodynamicLane.navierStokesForm}
Under the hypotheses of Theorem~\ref{cor:hydrodynamic-closure-law}, define
\[
\lambda_{\mathrm{ind}}
:=
\zeta_{\mathrm{ind}}-\frac{2}{d}\mu_{\mathrm{ind}}.
\]
Then the momentum law may be written in the standard compressible
Navier--Stokes form
\[
\partial_t(U_0 w)
+
\nabla\!\cdot(U_0 w\otimes w)
+
\nabla s
=
\nabla\!\cdot\Bigl(
\mu_{\mathrm{ind}}(\nabla w+\nabla w^\top)
+
\lambda_{\mathrm{ind}}(\nabla\!\cdot w)I
\Bigr)
-\nabla\!\cdot \mathsf R_{\mathrm{hyd}}.
\]
Writing
\[
\rho:=U_0,
\qquad
\rho u:=U_1,
\qquad
p:=s,
\]
this is the continuity-plus-momentum sector of compressible Navier--Stokes.
The familiar two-coefficient isotropic viscous stress is not inserted as an
extra constitutive hypothesis here: it is the forced trace/traceless endpoint
decomposition seen on the least-motion characteristic cut.
\end{corollary}

(Proof in Appendix~\ref{sec:appendix-dynamics}.)

\begin{corollary}[Exact Schur form of the selected hydrodynamic remainder]
\label{cor:hydrodynamic-schur-remainder-identity}\lean{HydrodynamicLane.schurRemainder}
Assume the hypotheses of Theorem~\ref{cor:hydrodynamic-closure-law}. Assume in
addition that the same selected datum satisfies a canonical selected Schur
bridge. Then there is a unique hydrodynamic-projected Schur correction
\[
\Sigma_{h,\mathrm{hyd}}(z)
\]
obtained by reading the common selected Schur correction \(\CCSigma{h}{z}\)
through the selected resolution-integrated hydrodynamic endpoint map, and
\[
\mathsf R_{\mathrm{hyd}}=\Sigma_{h,\mathrm{hyd}}(z).
\]
Accordingly,
\[
\Theta
=
-2\mu_{\mathrm{ind}}\,\mathsf D_0(w)
-\zeta_{\mathrm{ind}}(\nabla\!\cdot w)I
+
\Sigma_{h,\mathrm{hyd}}(z),
\]
and hence
\[
\partial_t(U_0 w)
+
\nabla\!\cdot(U_0 w\otimes w)
+
\nabla s
=
\nabla\!\cdot\Bigl(
2\mu_{\mathrm{ind}}\,\mathsf D_0(w)
+
\zeta_{\mathrm{ind}}(\nabla\!\cdot w)I
\Bigr)
-\nabla\!\cdot \Sigma_{h,\mathrm{hyd}}(z).
\]
\end{corollary}

(Proof in Appendix~\ref{sec:appendix-dynamics}.)
